\begin{table*}[!t]
\centering
\caption{Module~3 constants. Rate constants and pool sizes are paired: $k_{E,\mathrm{on}}[\mathrm{GEF_R}]_\mathrm{tot}$, $k_{R,\mathrm{on}}[\mathrm{GEF_R}]_\mathrm{tot}$ and $k_{I,\mathrm{on}}[\mathrm{GAP}]_\mathrm{tot}$ were held fixed when the pools were enlarged, so that only shot noise changed.}
\label{tab:m3const}
\footnotesize
\renewcommand{\arraystretch}{1.15}
\begin{tabularx}{\textwidth}{@{}l r l c >{\raggedright\arraybackslash}X@{}}
\toprule
Symbol & Value & Unit & Basis & Meaning, justification, source\\
\midrule
$[\mathrm{RasG}]_\mathrm{tot}$ & \num{38500} & $\mathrm{molec/voxel}$ & E & $5\times10^5$ per cell; large for low shot noise. The copy number is not well measured. \\
$[\mathrm{GEF_R}]_\mathrm{tot}$ & \num{138600} & $\mathrm{molec/voxel}$ & E & Excitor pool ($1.8\times10^6$ per cell). \mkt{M3-gefr}{RasGEFR activates RasG}~\cite{kae2007}. \\
$[\mathrm{GAP}]_\mathrm{tot}$ & \num{46200} & $\mathrm{molec/voxel}$ & E & Inhibitor pool ($6\times10^5$ per cell); \mkt{M3-nf1}{a RasGAP (NF1) is the global inhibitor}~\cite{zhang2008}. \\
$k_{E,\mathrm{on}}$ & \num{3.8e-3} & $\mu\mathrm{M}^{-1}\mathrm{s}^{-1}$ & C & G$\beta\gamma$-driven GEF activation; \mkt{M3-kort}{the initial Ras response requires G$\beta$}~\cite{kortholt2013}. \\
$k_{E,\mathrm{off}}$ & \num{0.5} & $\mathrm{s^{-1}}$ & C & GEF inactivation, $\tau_E=2$\,s: \mkt{M3-peak-fast}{the excitor is fast, so the Ras peak occurs within seconds}~\cite{kae2004,takeda2012}; close to \mkt{M3-tau-fit}{the fitted $0.4\,\mathrm{s^{-1}}$}~\cite{takeda2012}. \\
$k_{R,\mathrm{on}}$ & \num{0.333} & $\mu\mathrm{M}^{-1}\mathrm{s}^{-1}$ & C & GEF-catalysed GDP$\to$GTP exchange on RasG. \\
$k_{I,\mathrm{on}}$ & \num{7.67e-4} & $\mu\mathrm{M}^{-1}\mathrm{s}^{-1}$ & D & G$\beta\gamma$-driven GAP activation. It reads the \emph{same} input as 3.1, as \mkt{M3-prop}{an incoherent feed-forward loop requires}~\cite{takeda2012,goentoro2009}. Calibrated value $2.53\times10^{-4}$ multiplied by $3.03$ together with $k_{I,\mathrm{off}}$, which keeps the activated fraction; \mkt{M3-tau-pair}{the fit tied activation and deactivation rates in the same way}~\cite{takeda2012}. \\
$k_{I,\mathrm{off}}$ & \num{0.1} & $\mathrm{s^{-1}}$ & L & GAP inactivation, $\tau_I=10$\,s, \mkt{M3-tau-fit}{the fitted value, which sets the time of the return to the basal amount}~\cite{takeda2012} (until 2026-10: $0.033\,\mathrm{s^{-1}}$); \mkt{M3-slower-inh}{the inhibitor is slower than the excitor, which produces adaptation}~\cite{takeda2012}. Not yet recalibrated against the downstream modules. \\
$k_{I,\mathrm{b}}$ & \num{1.52e-3} & $\mathrm{s^{-1}}$ & D & Stimulus-independent GAP activity. Without it the resting Ras/PIP3 loop latches; consistent with \mkt{M3-basal-gap}{a basal RasGAP}~\cite{zhang2008,xu2021c2gap}. Calibrated as the minimum value giving the low rest state ($5\times10^{-4}$ at $\tau_I=30$\,s), multiplied by $3.03$ with $k_{I,\mathrm{off}}$, which keeps the basal activated fraction. \\
$K_\mathrm{NF1}$ & \num{0.2} & $\mu\mathrm{M}$ & L & \mkt{M3-nf1kin}{Michaelis constant of the NF1 GAP domain, $0.3\,\mu$M}~\cite{wiesmuller1992,coyle2016}; the value is rounded to $0.2$. Mammalian catalytic domain; no \emph{Dictyostelium} kinetics exist. \\
$k_\mathrm{cat}$ & \num{15.8} & $\mathrm{s^{-1}}$ & D & Derived from the resting distance to the zero-order switch, $\theta_\mathrm{rest}=0.5$. Only the product $k_\mathrm{cat}[\mathrm{GAP}^*]$ enters; it is 11$\times$ the mammalian \mkt{M3-nf1kin}{$k_\mathrm{cat}=1.4$\,s$^{-1}$}~\cite{coyle2016,wiesmuller1992}, so the active-GAP pool stands for more enzyme than is measured. \\
$k_\mathrm{PIP3,R}$ & \num{0.04} & $\mu\mathrm{M}^{-1}\mathrm{s}^{-1}$ & C & PIP3$\to$RasG feedback (3.3b), \mkt{M3-pip3-gef}{the PIP3 term of the Ras GEF reaction}~\cite{fukushima2019}. The value committed for this form in an earlier calibration (higher gains raised the variance but not the contrast, see text); not re-tuned after the changes of 2026-10. Kept small, since \mkt{M3-ly-red}{PI3K inhibition reduces but does not abolish Ras activation}~\cite{sasaki2004}. Expressed relative to the legacy PIP2 pool (Sec.~\ref{sec:m4}). \\
$k_{B,\mathrm{on}}$ & \num{0.0753} & $\mu\mathrm{M}^{-1}\mathrm{s}^{-1}$ & E & Brake activation by RasG-GTP (3.3c$_\mathrm{R}$); half-activation at 6400 RasG-GTP molecules/voxel, the driven level recorded in the code (design anchor, not tuned). Proposed brake. \\
$k_{B,\mathrm{off}}$ & \num{0.1} & $\mathrm{s^{-1}}$ & D & Brake decay ($\tau=10$\,s), from \mkt{M3-refr}{the recovery half-time of about 7\,s after a refractory period}~\cite{artemenko2016} as $\tau=t_{1/2}/\ln2$, assuming that B$^*$ governs the recovery; measured with an F-actin reporter after mechanical stimulation, not for Ras. \\
$k_{B,\mathrm{hyd}}$ & \num{39.8} & $\mu\mathrm{M}^{-1}\mathrm{s}^{-1}$ & D & Chosen so that at half activation the brake removes RasG-GTP at $0.6\times$ the rate of the LEGI GAP arm at the driven state (design target, not tuned). Proposed brake. \\
$[\mathrm{B}]_\mathrm{tot}$ & \num{1540} & $\mathrm{molec/voxel}$ & E & Brake pool; the mean braking action is independent of pool size, so it was raised to reduce noise. \\
$D_{\mathrm{GEF_R}},D_{\mathrm{RasG}}$ & \num{0.1} & $\mu\mathrm{m^2/s}$ & E & Local excitor and membrane-bound Ras (prenylated anchor). \\
$D_{\mathrm{GAP}},D_\mathrm{B}$ & \num{20} & $\mu\mathrm{m^2/s}$ & E & Global inhibitor and global brake. Cytosolic value; for B$^*$, a membrane-attached species in the code, it contradicts the local purpose of the brake (open, see 3.3c$_\mathrm{R}$--e); $0.4\,\mu\mathrm{m^2/s}$ gives the local brake. \\
\bottomrule
\end{tabularx}
\end{table*}
